The transition toward sixth-generation (6G) wireless networks is characterized by a paradigm shift from cell-centric to user-centric architectures, specifically through the deployment of cell-free massive multiple-input multiple-output (MIMO) systems \cite{bashar2019cell}. By distributing a high density of access points (APs) across a geographic area to cooperatively serve a smaller set of users, cell-free MIMO enhances spectral efficiency and provides more uniform coverage. This distributed architecture increasingly aligns with the Open Radio Access Network (O-RAN) framework, which advocates for the disaggregation of hardware and software through standardized interfaces, such as fronthaul and midhaul links \cite{femenias2020fronthaul}. However, the practical implementation of these systems is constrained by the transport network, where the high-volume data exchange required for centralized signal processing often exceeds the available capacity of the backhaul links \cite{kassab2019non}.

In practical cell-free deployments, backhaul and fronthaul bandwidths are inherently dynamic and finite, fluctuating according to network traffic and physical link conditions. Centralized processing typically necessitates the transmission of high-resolution raw or locally processed signals from APs to a central processing unit (CPU), creating a significant communication bottleneck \cite{liu2015joint}. Conventional quantization methods, such as fixed-rate uniform scalar quantization, are often suboptimal in this context as they fail to adapt to real-time channel state information (CSI) or the heterogeneous significance of individual AP observations \cite{park2013multihop}. Under stringent backhaul capacity limits, these static approaches result in severe quantization noise or the loss of critical information, thereby degrading the performance of multi-user detection (MUD) \cite{zhou2016fronthaul}.

Recent advancements in deep learning (DL) have provided robust frameworks for addressing the complexities of MIMO detection and signal compression. Model-driven DL architectures, such as those proposed in \cite{he2020model} and \cite{wei2020learned}, have demonstrated performance improvements over traditional iterative detectors by unfolding optimization algorithms into neural networks. Concurrently, neural compression and task-oriented communication frameworks have emerged to optimize data transmission for specific downstream objectives rather than perfect reconstruction \cite{shao2021task, bian2025towards}. For instance, variational information bottleneck (VIB) principles have been applied to extract task-relevant features while minimizing bit rates \cite{charvin2023exact}. Furthermore, progressive and successive refinement quantization schemes have been explored to adapt to varying link budgets \cite{sohrabi2022learning, fishel2025remote}. Despite these developments, a gap remains in the joint optimization of inference and resource allocation for cell-free MIMO. Existing works often decouple the quantization process from the spatial aggregation logic or rely on static bit-width assignments that do not fully exploit the spatial diversity of distributed APs under dynamic constraints \cite{qiao2024meta}.

To address these challenges, this paper introduces IB-CellFree-GNN, a task-oriented framework that jointly optimizes bit-width allocation and multi-user detection in backhaul-constrained cell-free MIMO systems. The framework prioritizes the transmission of information critical for the demodulation task, maximizing performance under fluctuating bandwidth limits. The primary contributions of this paper are as follows:
\begin{itemize}
    \item We develop a task-oriented successive refinement quantization scheme based on the VIB principle. This scheme decomposes local AP observations into hierarchical bit layers ordered by their relevance to the detection task, enabling flexible bit-slicing according to real-time capacity.
    \item We design a differentiable bit-width policy network that utilizes the Gumbel-Softmax reparameterization trick \cite{lyu2025joint}. This allows the CPU to learn an optimal, discrete bit-depth selection policy for each AP during training, effectively balancing detection accuracy against backhaul overhead.
    \item We introduce a bit-aware heterogeneous graph neural network (GNN) for robust signal aggregation. By incorporating reported bit-depths into the spatial attention mechanism, the GNN adaptively modulates the weights of incoming signals, allowing high-precision APs to compensate for the quantization noise of low-precision or failed links.
    \item Numerical evaluations indicate that the proposed framework achieves near-full-precision performance with an average bit-width of 2.5 bits. In high-interference scenarios, the method achieves a 67.3\% reduction in bit error rate (BER) compared to fixed 2-bit quantization and maintains robustness even when 50\% of the APs experience backhaul failure.
\end{itemize}

The remainder of this paper is organized as follows. Section II defines the system model and the backhaul-constrained detection problem. Section III details the proposed IB-CellFree-GNN framework, including the successive refinement quantizer and the policy network. Section IV presents the numerical results and performance analysis. Finally, Section V concludes the paper.